\section{Concentration of boundary-weighted moments}\label{sec:boundary}
We now pass from unweighted spatial traces to the boundary pairing for
the opposite-field chain. We will control two quantities together: the
fraction of an even boundary moment outside the leading spatial line,
and the defect of the physical Gibbs purity. The boundary coefficients
depend on temperature. We compare the weighted moments across temperatures
through physical partition-function identities, while contracting the
residual spatial powers at a fixed temperature.

Use the simple positive leading eigenvalue $\lambda_j$ and residual
ratio $\rho_j$ from Section~\ref{sec:periodic}. Let $P_j$ be the
orthogonal projection onto its line, and put
\begin{equation}\label{be:definitions}
 \begin{split}
 c_j&=\norm{P_jv_{b_j}}^2,\\
 B_k(j)&=\langle v_{b_j},(|X_{b_j}|/\lambda_j)^k v_{b_j}\rangle,
 \qquad e_k(j)=\frac{B_k(j)-c_j}{B_k(j)},\\
 A(m,j)&=\frac{Y_m(b_{j+1})}{Y_m(b_j)^2}.
 \end{split}
\end{equation}
We use $k,m$ even and at least two. Then
$B_k(j)=Y_k(b_j)/\lambda_j^k>0$, so $e_k(j)$ is defined even if
$c_j=0$. Also $0\le e_k(j)\le1$ and $0<A(m,j)\le1$.
Thus $e_k(j)$ is the residual fraction of a spatial boundary moment,
whereas $A(m,j)$ is the physical Gibbs purity of the opposite-field
chain at inverse temperature $b_j$.

The normalized absolute spatial eigenvalues are at most one. Thus
$B_m(j)\le B_k(j)$ when $m\ge k$. More precisely, if
$D_k(j)=B_k(j)-c_j$, the spectral theorem gives
\begin{equation}\label{be:residual}
 0\le D_m(j)\le\rho_j^{m-k}D_k(j)\qquad(m\ge k).
\end{equation}
The powers act on the residual spectrum; no upper bound for the total
boundary norm is needed in this comparison.

\begin{lemma}[Transfer between temperatures]\label{be:recurrence}
Let $j\ge0$ and let $m>k\ge2$ be even. Suppose $e_k(j)<1$ and
$0<q<1$ with $q\ge\rho_{j+1}^{m-k}$. Define, for
$0\le e<1$ and $0<t\le1$,
\begin{equation}\label{be:F}
 F(q,e,t)=\frac{q}{1-q}\left(\frac{1}{(1-e)^2t}-1\right).
\end{equation}
Then
\begin{equation}\label{be:step}
 e_m(j+1)\le F(q,e_k(j),A(m,j)).
\end{equation}
In particular, an upper bound strictly below one on the right proves
$c_{j+1}>0$ without presupposing it.
\end{lemma}
\begin{proof}
At the new temperature, \eqref{be:residual} yields
$B_m(j+1)-c_{j+1}\le q[B_k(j+1)-c_{j+1}]$.
Subtract and rearrange to obtain
\[
 e_m(j+1)\le\frac{q}{1-q}
                   \left(\frac{B_k(j+1)}{B_m(j+1)}-1\right).
\]
Write $r_j=\lambda_{j+1}/\lambda_j^2$. From the definitions,
\begin{equation}\label{be:ratio}
 \frac{B_k(j+1)}{B_m(j+1)}
 =\frac{A(k,j)}{A(m,j)}
   \left(\frac{B_k(j)}{B_m(j)}\right)^2 r_j^{m-k}.
\end{equation}
Section~\ref{sec:periodic} gives $r_j\le1$, and physical purity gives
$A(k,j)\le1$. The hypothesis $e_k(j)<1$ implies $c_j>0$, and
$B_m(j)\ge c_j$ gives
$B_k(j)/B_m(j)\le(1-e_k(j))^{-1}$.
Substitution proves \eqref{be:step}. Only the old overlap $c_j$ has
been divided by. The new overlap follows from
$e_m(j+1)<1$ and $B_m(j+1)>0$.
\end{proof}

To couple this residual estimate to physical purity, we use a shorter
reference length for the boundary defect. Recall $M_j=240\,2^{j-2}$ and
$u_j$ from Proposition~\ref{per:concentration}. For $j\ge2$ let
\begin{equation}\label{be:k}
 k_2=142,\qquad k_{j+1}=M_j.
\end{equation}
We track upper bounds $E_j$ on $e_{k_j}(j)$ and $\eta_j$ on
$1-A(M_j,j)$. The difference $M_j-k_j$ supplies residual contraction
at the next temperature. Its output at length $k_{j+1}=M_j$ then
controls the ratio $B_{2M_j}(j+1)/B_{M_j}(j+1)$ needed to improve
physical purity. The notation $\eta_j$ denotes only this physical
purity defect.

\begin{proposition}[Coupled boundary update]\label{be:coupled}
Suppose $j\ge2$, $0\le E_j,\eta_j<1$, and
\[
 e_{k_j}(j)\le E_j,\qquad A(M_j,j)\ge1-\eta_j.
\]
Let $E_{j+1},\eta_{j+1}\in[0,1)$ satisfy
\begin{align}
 E_{j+1}&\ge
 F\left(d(u_{j+1})^{(M_j-k_j)/M_j},E_j,1-\eta_j\right),
                                                   \label{be:Eupdate}\\
 \eta_{j+1}&\ge f\left(1-(1-\eta_j)(1-u_j)(1-u_{j+1})
                                      (1-E_{j+1})\right).
                                                   \label{be:etaupdate}
\end{align}
Then $e_{k_{j+1}}(j+1)\le E_{j+1}$ and
$A(M_{j+1},j+1)\ge1-\eta_{j+1}$.
\end{proposition}
\begin{proof}
Here $M_j>k_j$ and $n_{j+1}=M_j$. The periodic residual bound gives
\[
 \rho_{j+1}^{M_j-k_j}\le d(u_{j+1})^{(M_j-k_j)/M_j}<1.
\]
The function $F$ increases in $q,e$ and decreases in $t$ on the stated
domains. Lemma~\ref{be:recurrence} therefore proves
\eqref{be:Eupdate}'s conclusion, including $c_{j+1}>0$.

Put $M=M_j$. Exact cancellation in \eqref{be:definitions} gives
\begin{equation}\label{be:physical-cancellation}
 \frac{A(2M,j)}{A(M,j)}
 =r_j^M\frac{B_{2M}(j+1)}{B_M(j+1)}
             \left(\frac{B_M(j)}{B_{2M}(j)}\right)^2.
\end{equation}
The squared factor is at least one, and the first boundary ratio is
at least $c_{j+1}/B_M(j+1)\ge1-E_{j+1}$. Thus
Proposition~\ref{per:concentration} gives
\[
 A(2M,j)\ge(1-\eta_j)(1-u_j)(1-u_{j+1})(1-E_{j+1}).
\]
At the fixed physical length $2M$, passing from $b_j$ to $b_{j+1}$
normalizes the squares of the Gibbs probabilities.
Lemma~\ref{per:purity} bounds the resulting defect by $f$ of the
preceding defect. Since $M_{j+1}=2M$, this is
\eqref{be:etaupdate}.
\end{proof}
